\documentclass{article}
\usepackage{graphicx}
\usepackage{amsmath}
\usepackage[T2A]{fontenc}  % For Cyrillic fonts
\usepackage[utf8]{inputenc}  % For UTF-8 encoding
\usepackage[russian]{babel}  % For Russian language support
\usepackage{listings}
\usepackage{xcolor}
\usepackage{mathtools}

\lstset{
    language=Python,
    basicstyle=\ttfamily\small,
    keywordstyle=\color{blue},
    stringstyle=\color{red},
    commentstyle=\color{gray},
    morekeywords={self},
    frame=single,
    breaklines=True
}

\title{Algorithms 1 \\ Homework 2}
\author{Nikita Zagainov}
\date{September 2024}

\begin{document}

\maketitle
\section{Complexities, Fast Power, Fibonacci}
\begin{enumerate}
    \item
          \[
              13_{10} = 1101_{2}
          \]
          Применим алгоритм быстрого возведения в степень шаг за шагом, начиная с состояния $curr = 1$.
          \begin{enumerate}
              \item[1.] $b = 1$, $curr = curr^2 \cdot 7 \mod 167 = 7$
              \item[2.] $b = 1$, $curr = curr^2 \cdot 7 \mod 163 = 17$
              \item[3.] $b = 0$, $curr = curr^2 \mod 163 = 119$
              \item[4.] $b = 1$, $curr = curr^2 \cdot 7 \mod 163 = 143$
          \end{enumerate}
    \item Для поиска членов последовательности Фиббоначи для $k > 0$ от $F_0 = 0$ и $F_1 = 1$, можно пользоваться этой формулой:
          \[
              \begin{bmatrix}
                  F_{k+1} \\
                  F_{k}
              \end{bmatrix}
              =
              \begin{bmatrix}
                  1 & 1 \\
                  1 & 0
              \end{bmatrix}
              \begin{bmatrix}
                  F_{k} \\
                  F_{k-1}
              \end{bmatrix}
              =
              \begin{bmatrix}
                  1 & 1 \\
                  1 & 0
              \end{bmatrix}^k
              \begin{bmatrix}
                  1 \\
                  0
              \end{bmatrix}
          \]
          Для поиска членов при $k < 0$ справедливо:
          \[
              \begin{bmatrix}
                  F_{k} \\
                  F_{k-1}
              \end{bmatrix}
              =
              \begin{bmatrix}
                  1 & 1 \\
                  1 & 0
              \end{bmatrix}^{-1}
              \begin{bmatrix}
                  F_{k+1} \\
                  F_{k}
              \end{bmatrix}
              =
              \begin{bmatrix}
                  1 & 1 \\
                  1 & 0
              \end{bmatrix}^{k}
              \begin{bmatrix}
                  1 \\
                  0
              \end{bmatrix}
          \]
          То есть, даже для $k < 0$ данная формула применима

    \item Для решения задачи следует отсортировать котов по возрастанию массы.
          \begin{lstlisting}
    def insertion_sort(cats):
        for i in range(1, len(cats)):
            cur_mass = cats[i].mass
            j = i - 1
            while j >= 0 and cur_mass < cats[j].mass:
                cats[j + 1] = cats[j]
                j -= 1
            cats[j + 1] = key

        return cats
    \end{lstlisting}
          Этот алгоритм вернет массив из котов, отсортированных по возрастанию массы. Докажем, что он корректен математической индукцией. Предположим, что на каждой итерации первые $i$ элементов расположены в порядке возрастания. Очевидно, это правда для момента $i = 1$. Шаг индукции: предположим, что это свойство верно для $i = k$, тогда на переходе к $i = k + 1$ это свойство сохраняется, так как алгоритм берет $k + 1$ элемент и передвигает его в отсортированную часть массива, попарно сравнивая с каждым элементом по пути, пока элемент слева не окажется меньше. Таким образом, этот алгоритм корректен.

          Этот алгоритм имеет сложность $\Theta(n^2)$ где $n$ - количество котов, так как он имеет два вложенных цикла, итерирующих по всему массиву.

    \item
          \[
              f(n) = \sum_{k=0}^{n}{n \choose k} = \sum_{k=0}^{n}{n \choose k}1^{k}1^{n-k} = (1 + 1)^n = 2^n
          \]
    \item Пусть $g(n)$ - количество ``heh'', который напечатает вызов $f(n)$
          \begin{align*}
               & g(n) =                                                                                                                                                      \\
               & \sum_{i=1}^{\lfloor \log_2{n} \rfloor}{\sum_{j=0}^{2^{3i} - 1}{1}} + \sum_{i=1}^{\lfloor \frac{n}{2} \rfloor}{\sum_{j=0}^{(2i + 1)^2 - 1}{1}} =             \\
               & \sum_{i=1}^{\lfloor \log_2{n} \rfloor}{2^{3i}} + \sum_{i=1}^{\lfloor \frac{n}{2} \rfloor}{(2i + 1)^2} =                                                     \\
               & \frac{2^{3\lfloor \log_2{n} \rfloor} + 8}{2} \lfloor \log_2{n} \rfloor + \frac{(2\lfloor \frac{n}{2} \rfloor + 1) ^ 2 + 1}{2} \lfloor \frac{n}{2} \rfloor = \\
               & \Theta(n^3\log{n} + n^3) =                                                                                                                                  \\
               & \Theta(n^3\log{n})
          \end{align*}
\end{enumerate}

\end{document}
